% Handout for the M07 mini-project: the task, the requirements, the targets, and the
% two helpers printed from walk_kit.py so the paper and the file cannot disagree.
% The notebook carries none of this text.
% Targets are the Leiden means printed by tools/build_m07_mini_data.py.
% Compile with: xelatex -interaction=nonstopmode handout.tex   (or tools/build_worksheets.sh)
\documentclass[a4paper, 12pt]{extarticle}
\usepackage[top=0.8in, bottom=0.7in, left=0.85in, right=0.85in]{geometry}
\usepackage{fontspec}
\usepackage{listings}
\usepackage{xcolor}
\usepackage{enumitem}
\usepackage[hidelinks]{hyperref}

% Body face: Charter by file (XCharter), then the system fallbacks the other sheets use.
\IfFontExistsTF{XCharter-Roman.otf}{
  \setmainfont{XCharter}[
    Extension      = .otf,
    UprightFont    = *-Roman,
    BoldFont       = *-Bold,
    ItalicFont     = *-Italic,
    BoldItalicFont = *-BoldItalic]
}{
  \IfFontExistsTF{Charter}{\setmainfont{Charter}}{
    \IfFontExistsTF{Georgia}{\setmainfont{Georgia}}{}}}

% Code face, resolved once: Menlo on a Mac, otherwise the TeX default mono.
\IfFontExistsTF{Menlo}{\setmonofont{Menlo}[Scale=0.88]}{\setmonofont{Latin Modern Mono}[Scale=0.95]}

\definecolor{rule}{HTML}{D8D4CC}
\definecolor{muted}{HTML}{6B6B6B}
\definecolor{rust}{HTML}{B5482F}

\lstset{
  language=Python,
  basicstyle=\ttfamily\small,
  commentstyle=\color{muted},
  keywordstyle=\bfseries,
  stringstyle=\color{rust},
  showstringspaces=false,
  columns=fullflexible,
  keepspaces=true,
  frame=single,
  rulecolor=\color{rule},
  framesep=7pt,
  xleftmargin=7pt,
  xrightmargin=7pt,
  aboveskip=0.6em,
  belowskip=0.4em,
}

\setlength{\parindent}{0pt}
\setlength{\parskip}{0.5em}
\pagestyle{empty}

\begin{document}

{\Large\bfseries M07 Mini project: community detection with a random walk}

\textbf{Task.} Design a community detection algorithm that uses a random walk, and
test it on three networks, \texttt{net1}, \texttt{net2} and \texttt{net3}. You are
not told what they are; working that out from the data is part of the task. Each
comes with a ground-truth grouping of its nodes.

\textbf{Requirements}
\begin{enumerate}[leftmargin=1.6em, itemsep=0.2em, topsep=0.2em]
  \item The algorithm uses a random walk: a function that simulates walks, or the
        transition matrix $P$ of the walk.
  \item It is a community detection algorithm: from the edge table it returns one
        group label per node. The number of groups is not given.
  \item It is one algorithm. You may set parameters separately for each network, but
        only parameters that change how the walk behaves, for example its length, a
        restart probability, or the probability of stepping back. A parameter that
        switches between algorithms, or that sets the number of groups, is not allowed.
  \item No library community detection: no \texttt{community\_*} in igraph, no
        community functions in networkx, no \texttt{SpectralClustering}. Linear
        algebra, $k$-means, hierarchical clustering and mixture models are allowed.
  \item The algorithm does not read the ground truth.
\end{enumerate}

\textbf{Notebook.} It loads the three edge tables and the ground truth as pandas
tables. Write your code in the empty cell. Put the result in \texttt{answer}: a
dictionary from the network name to a list with one group label per node, in node
order. Run that cell; it prints one score per network.

\textbf{Score and targets.} Each network is scored by the normalized mutual
information (NMI) between your groups and the ground truth: 1 if they agree exactly,
near 0 if they share no information. The targets are the scores of modularity
maximisation (Leiden), averaged over 10 runs.

\begin{center}
\begin{tabular}{lcc}
  \hline
  network & target & \\
  \hline
  \texttt{net1} & 0.78 & \\
  \texttt{net2} & 0.00 & any score above 0 \\
  \texttt{net3} & 0.68 & \\
  \hline
\end{tabular}
\end{center}

Beating the targets is the goal. Submit the notebook through the Google Form linked
from the Brightspace assignment, whatever the scores are.

\newpage

\textbf{Two helpers.} Use them, change them, or ignore them. They work on an edge
table \texttt{edges} with columns \texttt{source} and \texttt{target}, for example
\texttt{edges["net1"]}. Both need only \texttt{numpy} and \texttt{scipy}.

\lstinputlisting[firstline=3]{walk_kit.py}

\textbf{Using them}
\begin{lstlisting}
A = to_adjacency(edges["net1"])  # sparse adjacency matrix, n x n
P = transition_matrix(A)         # P[i, j]: probability of moving from i to j
walk = random_walk(A, 0, 10)     # a list of 11 nodes, starting at node 0
\end{lstlisting}

\end{document}
